% TOP_PROBLEM: {"id": 3000068, "title": "Rainbow matchings in bipartite graphs", "category_id": 2, "category": {"id": 2, "name": "combinatorics", "display_name": "Combinatorics", "description": "Counting problems, graph theory, discrete structures.", "slug": "combinatorics", "order_index": 2, "created_at": "2026-07-31T15:26:25.671Z"}, "status": "open", "problem_number": "AMR-029-0068", "set_id": 13}
% FIRST_POSTED: 2026-10-01
% ENTRY_KIND: statement_only
% UnsolvedMath ID 3000068; source catalogue code AMR-029-0068
% !TeX program = lualatex
% Definition and sources only; this file is not an LLM solution attempt.
\documentclass[11pt,a4paper]{article}
\usepackage[margin=25mm]{geometry}
\usepackage{fontspec}
\setmainfont{lmroman10-regular.otf}[BoldFont=lmroman10-bold.otf,
 ItalicFont=lmroman10-italic.otf,BoldItalicFont=lmroman10-bolditalic.otf]
\usepackage{amsmath,amssymb,mathtools}
\usepackage{unicode-math}
\setmathfont{latinmodern-math.otf}
\AtBeginDocument{\let\setminus\smallsetminus}
\usepackage{xurl,hyperref}
\hypersetup{colorlinks=true,urlcolor=blue}
\setcounter{secnumdepth}{0}
\setlength{\parindent}{0pt}
\setlength{\parskip}{5pt}
\setlength{\emergencystretch}{3em}
\begin{document}
\section*{Rainbow matchings in bipartite graphs}\label{3000068}
{\small UnsolvedMath ID 3000068 · AMR-029-0068 · Combinatorics · AMR Open Problem Lists}
\subsection{Definitions and mathematical statement}
Let $G=(U\sqcup W,E)$ be a finite bipartite graph. A matching is
a set of edges with no shared endpoint. Fix an integer $k\ge1$
and suppose $M_1,\ldots,M_k$ are pairwise edge-disjoint matchings,
each with exactly $k+1$ edges.

The conjecture asks whether one can choose $e_i\in M_i$ for
$i=1,\ldots,k$ so that the chosen edges also form a matching:
\[
                         e_i\cap e_j=\varnothing
                                           \quad(i\ne j),
\]
where an edge is identified with its two endpoints in this condition.
Such a choice is a full rainbow matching, one edge of every colour
when edges in $M_i$ are coloured $i$.

Disjointness of the given matchings refers to their edge sets, not
to their vertex sets; different $M_i$ may meet the same vertices.
If they were vertex-disjoint the conclusion would be immediate.
Extra edges of the ambient graph outside their union have no role.
The size requirement can equivalently be at least $k+1$, since
each input matching may first be truncated. The requested result
is exact for every $k$, so a matching missing a few colours, or a
theorem requiring substantially larger colour classes, does not
settle it. Bipartiteness and the matching condition within each
colour class are retained hypotheses.
\subsection{Short English statement}
Do $k$ edge-disjoint matchings of size $k+1$ in a bipartite graph
always admit a matching choosing one edge from each?
\subsection{Sources}
\begin{itemize}
\item[S1] ULAM.AI and the UnsolvedMath Contributors, supplied problems.json, record 3000068 (AMR-029-0068), AMR Open Problem Lists. Catalogue identity and source transcription; CC BY 4.0.
\url{https://huggingface.co/datasets/ulamai/UnsolvedMath}.
\item[S2] Egres Open - List of open problems, Open-problem page: Rainbow matchings in bipartite graphs. Original source indexed by AMR-029-0068.
\url{https://oldlemon.cs.elte.hu/egres/open/List_of_open_problems}.
\item[S3] EGRES Open, Rainbow matchings in bipartite graphs. Individual source page and explanatory remarks.
\url{https://lemon.cs.elte.hu/egres/open/Rainbow_matchings_in_bipartite_graphs}.
\end{itemize}
\end{document}
